% Response to the review's methodological objection about chromosome-end
% restriction distorting acrocentrics and PARs. Drop into the response letter
% and/or the Discussion; not applied automatically.
% Numbers: submission/supporting_material/ED_community_tier_table.tsv

We agree that restricting analysis to chromosome ends requires that the
acrocentric short arms and the pseudoautosomal regions be treated separately,
and we have quantified why. The ratio of non-homologous to homologous divergence
separates the genome into three regimes (Extended Data Table~\ref{tab:community_tiers}).
In the acrocentric short arms that ratio is 1.00: copies on different
chromosomes differ by 1.135\%, while homologous copies of the same arm differ by
1.130\%. In PAR1 it is 0.99 and in PAR2 1.43. These are the regions in which
recombination between non-allelic copies is obligate rather than ectopic ---
acrocentric short arms exchange ribosomal DNA arrays, and the pseudoautosomal
regions cross over in every male meiosis --- and at that intensity the sequences
converge on the allelic floor of within-population diversity (0.10--0.16\%),
which no amount of exchange can lower further. The remaining arms form a
continuum from ratio 1.2 to 44, i.e. exchange at 2--85\% of the obligate rate.
The separation is therefore measured rather than assumed, and the acrocentric
and sex-chromosome arms are reported as a distinct regime instead of being
pooled with or excluded from the ectopic-exchange analysis.
